% Source printed pp. 28-31, entire Section 1.4.
\section{模群胚}\label{src:1.4}
\begin{epigraph}
我目前得出的实际结论是: 每当按照我的判别准则, 尽管具有良好的平坦性, 固有性以及可能的非奇异性假设, 却仍不存在用来分类某些结构(非奇异完备簇, 向量丛等)的(整体或无穷小)变化的模簇(更准确地说, 模概形)时, 原因仅仅在于这些结构具有自同构, 使下降技术无法奏效.
\epigraphauthor{Alexander Grothendieck, 致 Serre 的信, $1959$ 年~\cite[第 94 页]{CS01}}
\end{epigraph}
现在我们改变观点: 我们不只指定两个对象\textit{何时}被等同, 还要指定它们\textit{如何}被等同!

模空间最令人期待的性质之一, 是存在万有族(见 \S\bookref{1.3.4-subsection}); 而自同构的存在阻碍了万有族的存在(见 \S\bookref{1.3.5-subsection}). 在描述中纳入自同构, 将使我们能够克服这个问题. 定义一个模群胚, 需要明确: 
\begin{enumerate}[(1)]
\item 对象; 以及
\item 任意两个对象之间的等同所组成的集合(可能为空).
\end{enumerate}
不久, 我们将把上一节的函子观点与这一群胚观点结合起来, 定义模叠.

\subsection{群胚}\label{src:1.4.1-subsection}
用来编码对象及其等同的一种便利的数学结构, 就是群胚.
\begin{definition}\label{src:1.4.1}
\newterm{群胚}{groupoid}是一个范畴 $\mathcal C$, 其中每个态射都是同构.
\end{definition}
两个群胚 $\mathcal C_1$ 和 $\mathcal C_2$ \textit{等价}, 是指存在范畴等价 $F:\mathcal C_1\xrightarrow{\sim}\mathcal C_2$, 即存在函子 $G:\mathcal C_2\to\mathcal C_1$, 使 $F\circ G$ 和 $G\circ F$ 都同构于恒等函子; 等价地, $F$ 完全忠实且本质满.
\begin{example}[集合是群胚]\label{src:1.4.2}
若 $\Sigma$ 为集合, 则范畴 $\mathcal C_\Sigma$ 为群胚, 其对象是 $\Sigma$ 的元素, 而态射只有恒等态射. 若存在范畴等价 $\mathcal C\to\mathcal C_\Sigma$, 则称群胚 $\mathcal C$ \textit{等价于集合 $\Sigma$}.
\end{example}
\begin{example}[分类群胚]\label{src:1.4.3}
若 $G$ 为群, $G$ 的\newterm{分类群胚}{classifying groupoid} $BG$ 定义为只有一个对象 $\star$ 的范畴, 满足 $\operatorname{Aut}(\star)=\operatorname{Mor}(\star,\star)=G$.
\end{example}
\begin{example}\label{src:1.4.4}
由有限集组成, 以双射为态射的范畴 $\mathsf{FB}$ 是群胚. $\mathsf{FB}$ 的同构类与 $\mathbb N$ 一一对应.

而 $\operatorname{Aut}(\{1,\ldots,n\})=S_n$ 是置换群.
\end{example}
\begin{example}[射影空间]\label{src:1.4.5}
射影空间等同于过原点的直线 $L\subset\mathbb A^{n+1}$ 的模群胚, 其中态射只有恒等映射. 也可以将对象取为非零线性映射 $x=(x_0,\ldots,x_n):\mathbb C\to\mathbb C^{n+1}$, 且存在唯一态射 $x\to x'$ 当且仅当 $\operatorname{im}(x)=\operatorname{im}(x')\subseteq\mathbb C^{n+1}$(即存在 $\lambda\in\mathbb C^*$, 使 $x'=\lambda x$).
\end{example}

\subsection{轨道的模群胚}\label{src:1.4.2-subsection}
\begin{example}[轨道的模群胚]\label{src:1.4.6}
给定群 $G$ 在集合 $X$ 上的作用, 定义\textit{轨道的模群胚} $[X/G]$\footnote{我们用方括号区分群胚商 $[X/G]$ 与集合商 $X/G$. 后面, 当 $G$ 为代数群, $X$ 为概形时, $[X/G]$ 将表示商叠(它总是存在), 而 $X/G$ 将表示商空间(如果存在).}, 其对象为所有元素 $x\in X$, 并规定 $\operatorname{Mor}(x,x')=\{g\in G\mid x'=gx\}$.
\end{example}
\begin{figure}[htbp]
\centering\begin{tikzpicture}[x=1cm,y=1cm,>=Stealth,font=\small]
\begin{scope}
\draw (-1.2,1.4)..controls (.2,1.4) and (.2,.6)..(1.2,.6);
\draw (-1.2,.6)..controls (.2,.6) and (.2,1.4)..(1.2,1.4);
\node[right] at (1.35,1) {$\mathbb A^1$};
\draw[->] (0,.55)--(0,.13);
\draw (-1.2,0)--(1.2,0);
\fill (0,0) circle (1.8pt);
\node[below,blue] at (0,0) {$\mathbb Z/2$};
\node[right] at (1.35,0) {$[\mathbb A^1/(\mathbb Z/2)]$};
\end{scope}
\begin{scope}[xshift=7cm]
\draw (-1.2,1)--(1.2,1);
\fill (0,1) circle (1.8pt);\node[above] at (0,1) {$0$};
\node[right] at (1.35,1) {$\mathbb A^1$};
\draw[->] (0,.8)--(0,.23);
\fill (0,0) circle (1.8pt);
\shade[ball color=gray] (.7,0) circle (.10);
\draw[gray,densely dotted] (.6,0)--(.1,0);
\node[below,blue] at (0,0) {$\mathbb G_m$};
\node[below,blue] at (.7,0) {$\{1\}$};
\node[right] at (1.35,0) {$[\mathbb A^1/\mathbb G_m]$};
\end{scope}
\end{tikzpicture}

\caption{$\mathbb Z/2=\{\pm1\}$ 和 $\mathbb G_m$ 在 $\mathbb C$ 上的 $\mathbb A^1$ 上的数乘作用, 自同构群以蓝色列出. 注意, $[\mathbb A^1/\mathbb G_m]$ 有两个对象同构类, 即 $0$ 和 $1$, 对应于两个轨道, 即 $0$ 和 $\mathbb A^1\setminus0$. 若赋予集合 $\mathbb A^1/\mathbb G_m$ 商拓扑, 则 $0\in\overline{\{1\}}$.}\label{src:1.4.7}
\end{figure}
\begin{exercise}[简单]\label{src:1.4.8}
证明例~\bookref{1.4.6}~中轨道的模群胚 $[X/G]$ 等价于一个集合, 当且仅当 $G$ 在 $X$ 上的作用自由.
\end{exercise}
\begin{example}\label{src:1.4.9}
考虑具有两个对象 $x_1$ 和 $x_2$ 的范畴 $\mathcal C$, 对 $i,j=1,2$ 有 $\operatorname{Mor}(x_i,x_j)=\{\pm1\}$, 态射的复合由乘法给出. 那么 $\mathcal C$ 等价于 $B(\mathbb Z/2)$.
\end{example}
\begin{figure}[htbp]
\centering\begin{tikzpicture}[>=Stealth,font=\small]
\node[circle,fill,inner sep=3pt,label=below:$x_1$] (a) at (0,0) {};
\node[circle,fill,inner sep=3pt,label=below:$x_2$] (b) at (3,0) {};
\node[circle,fill,inner sep=3pt,label=below:$x$] (c) at (7.1,0) {};
\foreach \v in {a,b,c} {
\path[->] (\v) edge[out=150,in=30,looseness=16] node[above] {$1$} (\v);
\path[->] (\v) edge[out=135,in=45,looseness=10] node[below,pos=.5] {$-1$} (\v);
}
\path[->] (a) edge[bend left=18] node[above] {$1$} (b);
\path[->] (a) edge[bend left=32] node[above] {$-1$} (b);
\path[->] (b) edge[bend left=18] node[below] {$1$} (a);
\path[->] (b) edge[bend left=32] node[below] {$-1$} (a);
\draw[->] (4,.15)--(5.9,.15) node[midway,above] {$\sim$};
\node at (4.95,-.35) {$x_i\longmapsto x$};
\end{tikzpicture}

\caption{群胚的一个等价.}\label{src:1.4.10}
\end{figure}
\begin{exercise}[简单]\label{src:1.4.11}
在例~\bookref{1.4.9}~中, 证明 $\mathcal C$ 与 $[(\mathbb Z/2)/(\mathbb Z/4)]$ 或 $[(\mathbb Z/2)/(\mathbb Z/2\times\mathbb Z/2)]$ 之间存在一个在对象上诱导双射的范畴等价, 其中作用由满射 $\mathbb Z/4\to\mathbb Z/2$ 或 $\mathbb Z/2\times\mathbb Z/2\to\mathbb Z/2$ 给出.
\end{exercise}
\begin{example}[射影空间作为商]\label{src:1.4.12}
射影空间的模群胚(例~\bookref{1.4.5})也可以描述为轨道的模群胚 $[(\mathbb A^{n+1}\setminus0)/\mathbb G_m]$. 还可以考虑商群胚 $[\mathbb A^{n+1}/\mathbb G_m]$, 它等价于如下群胚: 对象为(可能为零的)线性映射 $x=(x_0,\ldots,x_n):\mathbb C\to\mathbb C^{n+1}$, 且 $\operatorname{Mor}(x,x')=\{t\in\mathbb C^*\mid\text{对所有 }i\text{ 有 }x'_i=tx_i\}$. 由此, $\mathbb P^n$ 是 $[\mathbb A^{n+1}/\mathbb G_m]$ 的一个子群胚.
\end{example}
\begin{exercise}[简单, 适于练习]\label{src:1.4.13}
若群 $G$ 在集合 $X$ 上作用, $x\in X$ 为稳定子是 $G_x$ 的点, 证明存在完全忠实函子 $BG_x\to[X/G]$. 若作用传递, 证明该函子为等价.
\end{exercise}
按定义, 群胚的\textit{态射} $\mathcal C_1\to\mathcal C_2$ 是一个函子. 范畴 $\operatorname{Mor}(\mathcal C_1,\mathcal C_2)$ 的对象为函子, 态射为自然变换.
\begin{exercise}[简单]\label{src:1.4.14}
若 $\mathcal C_1$ 和 $\mathcal C_2$ 为群胚, 证明 $\operatorname{Mor}(\mathcal C_1,\mathcal C_2)$ 为群胚.
\end{exercise}
\begin{exercise}[中等, 适于练习]\label{src:1.4.15}
若 $H$ 和 $G$ 为群, 证明存在等价
\[
\operatorname{Mor}(BH,BG)=
\coprod_{\phi\in\operatorname{Hom}(H,G)/G} B\bigl(C^G(\operatorname{im}\phi)\bigr),
\]
其中, $\operatorname{Hom}(H,G)/G$ 表示同态 $H\to G$ 模去 $G$ 的共轭作用后的等价类, $C^G(\operatorname{im}\phi)$ 表示 $\operatorname{im}\phi$ 在 $G$ 中的中心化子.
\end{exercise}
\begin{exercise}\label{src:1.4.16}
给出群 $H$ 和 $G$ 分别在集合 $X$ 和 $Y$ 上作用的一个例子, 以及群胚映射 $[X/H]\to[Y/G]$, 使该映射不能由群同态 $\phi:H\to G$ 与一个 $\phi$-等变映射 $X\to Y$ 产生.
\end{exercise}

\subsection{模群胚的例子}\label{src:1.4.3-subsection}
\begin{example}[光滑曲线的模群胚]\label{src:1.4.17}
对象为 $\mathbb C$ 上亏格为 $g$ 的光滑, 连通射影曲线. 对两条曲线 $C,C'$, 态射集合定义为同构集合
\[
\operatorname{Mor}(C,C')=\{\text{同构 }\alpha:C\xrightarrow{\sim}C'\}.
\]
\end{example}
\begin{example}[曲线上向量丛的模群胚]\label{src:1.4.18}
对象为固定曲线 $C$ 上秩为 $r$, 次数为 $d$ 的向量丛 $E$, 态射为向量丛的同构.
\end{example}
\begin{example}[无标号三角形的模群胚]\label{src:1.4.19}
回到例~\bookref{1.2.12}~中相似关系下无标号三角形的模 $M^{\mathrm{unl}}$. 回忆我们已经引入了无标号三角形的族, 并证明了这个函子不可表(例~\bookref{1.3.6}).

定义相似关系下\textit{无标号三角形的模群胚}, 记作 $\mathcal M^{\mathrm{unl}}$(注意花体), 其中对象为无标号三角形, 态射为相似变换. 例如, 等腰三角形和等边三角形的自同构群分别为 $\mathbb Z/2$ 和 $S_3$. 我们可以画出与图~\bookref{1.2.14}~基本相同的图, 只是还要记录自同构.
\end{example}
\begin{figure}[htbp]
\centering\begin{tikzpicture}[scale=5.5]
\fill[gray!30] (0,1)--(.5,.5)--(2/3,2/3)--cycle;
\draw[->] (0,0)--(1.08,0) node[right] {$a$};
\draw[->] (0,0)--(0,1.1) node[above] {$b$};
\draw[densely dotted] (0,1)--(.5,.5);
\draw[blue!35,line width=2pt] (0,1)--(2/3,2/3)--(.5,.5);
\draw[magenta,thick,domain=0:{2-sqrt(2)},samples=100] plot (\x,{(2-2*\x)/(2-\x)});
\fill[green!75!black] (2/3,2/3) circle (.012);
\foreach \v/\t in {0.5/{1/2},0.6666667/{2/3}} {
\draw (\v,0)--(\v,-.02) node[below] {$\t$};
\draw (0,\v)--(-.02,\v) node[left] {$\t$};
}
\node[left] at (0,1) {$1$};
\node[font=\small,rotate=-45,fill=white,inner sep=1pt] at (.2,.75) {$a+b=c$};
\node[font=\small,rotate=-27,above] at (.25,.9) {$b=c$};
\node[font=\small,rotate=45,below] at (.59,.57) {$a=b$};
\node[blue,rotate=-27,above] at (.44,.82) {$\mathbb Z/2$};
\node[blue,rotate=45,below] at (.66,.58) {$\mathbb Z/2$};
\node[right,green!75!black] at (.69,.68) {$S_3$};
\draw[->,gray] (.93,.86)--(.68,.69);\node[right,font=\small] at (.8,.89) {等边};
\draw[->] (.16,.25)--(.43,{(2-2*.43)/(2-.43)});
\node[font=\small] at (.17,.18) {$a^2+b^2=c^2$};
\end{tikzpicture}

\caption{模群胚 $\mathcal M^{\mathrm{unl}}$ 的示意图, 标出了非平凡自同构群.}\label{src:1.4.20}
\end{figure}
存在从模群胚到模集合的函子
\[
\mathcal M^{\mathrm{unl}}\longrightarrow M^{\mathrm{unl}},
\]
它在对象的同构类上给出等价, 并将所有态射压缩为恒等态射. 这是粗模空间的第一个例子.
\begin{exercise}\label{src:1.4.21}
回忆~\eqref{src:1.2.6}~中相似关系下带标号三角形的模集合 $M^{\mathrm{lab}}$ 的描述, 证明 $S_3$ 在相似关系下带标号三角形的模集合 $M^{\mathrm{lab}}$ 上有自然作用, 并且存在群胚的等同 $\mathcal M^{\mathrm{unl}}\cong[M^{\mathrm{lab}}/S_3]$.
\end{exercise}
\begin{exercise}\label{src:1.4.22}
定义有向三角形的模群胚, 并研究它与带标号及无标号三角形的模群胚之间的关系.
\end{exercise}
关于三角形的模叠的更详细阐述, 见~\cite{Beh14}.
